% Extracted missing sections from Chapter 3 (Deciphering the Inscriptions)
% Auto-generated from OCR cleaned files (Phase 3)
% Extracted sections: 3.3, 3.3.2, 3.6, 3.6.1, 3.7, 3.7.1.1, 3.7.3
%

% source: scan p. 76, section 3.3
\section[Reference Works]{\origsecnum{3.3}Reference Works}
\label{sec:chapters-ch03:3.3}

OverKaizuka (1967, p. 204, estimates that we know about 3,000 graphs on the oracle bones. about 800 of which can be interpreted; Chang Ping-ch'üan (1967), p. 829, refers to 3,320 graphs, of which we can recognize about 1,000, with 500 to 700 more identifiable but in dispute. Hu Houhsüan (1945), p. 5b, speculated, extravagantly in my view, that the full oracle-bone vocabulary, were we able to recover it, would have included 5,000 to 6,000 words, of which we would be able to recognize from 2,000 to over 3,000. No exactitude is possible in such counts (for other examples, see Chang Ch'un-shu [1972], p. 283); nor is anything more than a general estimate necessary. See too n. 20. three thousand Shang graphs have been identified on the oracle bones, and
about half of these may be translated or identified with varying degrees of certainty. The
dictionaries and concordances listed below are indispensable to the study of the inscriptions.

% source: scan p. 78-79, section 3.3.2
\subsection[Concordances and Compendiums]{\origsecnum{3.3.2}Concordances and Compendiums}
\label{sec:chapters-ch03:3.3.2}

Inkyo bokuji sōrui (abbreviated S) is a concordance of the oracle inscriptions in
sixty-three collections. 19 Over three thousand words have been classified. 20 Under each

entry, Shima Kunio has hand copied, in a standardized oracle-bone script (examples of
which grace the margins of this book) all the inscriptions in which the graph in question
appears, giving the citation for each. There are inevitable errors, inconsistent transcriptions, and omissions; some entries for the more common graphs have been abbreviated
(this is noted); and inscriptions published since 1959, together with some minor collections
published only in journals, are not included.It is true that the second edition identifies rejoined Yi-pien fragments that appear in the Ping-pien volumes published down to 1967, i.e., down to Ping-pien 512 (for indexes to these rejoinings, see n. 72); but in terms of new fragments (as opposed to new pieces; cf. appendix 3, n. 1) no collection published since 1959 is included. 21 But the work's fundamental contribution
is threefold. First, it attempts to provide a comprehensive listing of word usage in context.
Second, it provides control over word frequency (except, as already noted, in the case of
common graphs). Third. since Shima copies the inscriptions under each entry in chronological order (cf. sec. 4.1.1), it provides (within the limits imposed by the accuracy of his
transcription) a rapid introduction to the epigraphic evolution, if any, of individual
graphs (cf. sec. 4.3.1.4 and to their periodicity (appendix 5).
The work is also of major importance because, unlike the oracle-bone dictionaries
described above (sec. 3.3.1), it allows a graph entry to be found, even when the seal or modern
form of the graph is not known. This may be done by using the rear index (pp. 590—601),
arranged under the 164 oracle-bone ``radicals'' devised by Shima. Once the entry has
been found, researchers may study all occurrences of the graph in context; they may move
on to the commentaries which have been published for certain of the inscriptions cited;
or they may, in the 840 cases for which the cross reference is provided, move to the relevant
page in Li Hsiao-ting's Chia-ku wen-tzu chi-shih.Examples of how this works are given in secs. 3.6.1 to 3.6.3; see too n. 12. For my fuller description and evaluation of Inkyo bokuji sōrui, see Keightley (1969a) and RBS 12 (forthcoming).22 Inkyo bokuji sōrui and Chia-ku wen-tzu
chi-shih are complementary reference works. Neither should be relied on as a primary
source for the inscriptions or commentaries. But they are invaluable guides to the sources
and scholarship, the two fundamental works in the library of any oracle-bone scholar.
Based on a corpus of fifty-eight published and seven unpublished collections, Yin-tai
chen-pu jen-wu t'ung-k'aoJao Tsung-yi Yin-tai chen-pu jen-wu t'ung-k'ao, 2 vols. (Hongkong, 1959). 23 prints in modern transcription all the divination charges associated
with the names of 137 diviners.Cf. ch. 2, n. 13. 24 Within each diviner's section, the inscriptions are subdivided by topic: rain, the ten-day week, eclipses, hunting, and so on. A compendium
rather than a concordance, Jao Tsung-yi's work is less useful than Shima's; first, because
it is organized by diviner and topic rather than by words; second, because divinations
lacking a diviner's name are not included; third, because his modern-character transcriptions, useful though they are to the beginner, introduce a considerable element of interpretation. But it is a valuable supplement, especially for studying the role of particular
diviners. Once the diviner's name and the topic of a particular charge are known, a modern
transcription may generally be found without difficulty.For a more detailed assessment of Yin-tai chen-pu jen-wu t'ung-k'ao, see Noel Barnard in Monumenta Serica 19 (1960), pp. 480-486. _ 25

% source: scan p. 88, section 3.6
\section[How to Read the Graphs on a Bone or Shell Fragment]{\origsecnum{3.6}How to Read the Graphs on a Bone or Shell Fragment}
\label{sec:chapters-ch03:3.6}

Nothing appears more arcane, at first glance, than an oracle-bone rubbing. I
will try to demystify it in this section with a step-by-step analysis of how an actual turtle-shell piece, composed of two fragments (fig. 20), may be identified, deciphered, and placed
in context. The flow chart (fig. 21: summarizes the possible approaches.

% source: scan p. 88, section 3.6.1
\subsection[How to Identify the Fragment]{\origsecnum{3.6.1}How to Identify the Fragment}
\label{sec:chapters-ch03:3.6.1}

Orienting the piece is rarely a problem since rubbings are normally printed
correctly.Fragments have been printed upside down, especially in the earliest collections. See, for example, T'ieh-yün 23.4, 28.4, and the other cases listed in Yen Yi-p'ing's preface to the Yi-wen (Taipei, 1959) reprint of T'ieh-yün. Errors of this sort appear with some frequency in secondary works, e.g., Eisenberger (1938), p. 67 (the inscription is probably fake); Ch'en Chih-mai (1966), fig. 4a (whole plastron inverted); William S. Y. Wang (1972), p. 53 (bottom two fragments, left column, inverted). 64 If necessary, orientation is best done either by identifying certain graphs, or
graph elements, whose correct orientation is known, or by orienting the crack numbers
(near the top of the pu crack; sec. 2.4) and crack notations (near the bottom of the pu
crack; sec. 2.6) correctly. In figure 20, the crack number 6 and the shang-chi at the
left are readily placed. The pu cracks all run left; and the crack numbers 6, 7, and 8 run
right; these clues, together with the general configuration of the piece, lead us to suppose
that it comes from the right side of a shell in the region of the hyoplastron or hypoplastron
(fig.See ch. 2, n. 122. 3).65

% source: scan p. 93, section 3.7
\section[Reading Practices]{\origsecnum{3.7}Reading Practices}
\label{sec:chapters-ch03:3.7}

Translation not only involves placing inscriptions in the context of other inscriptions. It also involves understanding the role that divination and the carving of the inscriptions played in Shang culture; it involves understanding, as far as we can, what the diviner
thought he was doing.For a preliminary analysis of these issues, which will be considered more fully in Studies, see Keightley (1973a) and (1975c). Cf. too, Takashima (1976), who proposes a rationale for some of the charges in the set discussed below. 82 It is instructive in this regard to examine one set of inscriptions
in detail in an initial attempt to explore the rationale underlying Shang divination and to
indicate the kinds of questions that we need to ask.
I select the set of inscriptions carved on the five plastrons Ping-pien 12-21 because,
when the charges on the front are considered together with the subcharges (sec. 3.7.4.2)
and prognostications on the back, they form one of the most informative multiple-plastron
sets we now possess. (The fact that these five plastrons form a set [sec. 2.5] is demonstrated
by the similarity of the divination charges and the ordinal sequence of crack numbers
[sec. 2.4], which are all 1 on Ping-pien 12, all 2 on Ping-pien 14, and so on.)
The analysis which follows should, if possible, be studied with the Ping-pien rubbings
and Chang Ping-chüan's k'ao-shih at one's side; I do not always follow his interpretation,
and I make no attempt to reproduce his extensive notes about the identity of the various
figures and statelets involved. My concerns are primarily with the divination process.
3.7.1

% source: scan p. 93, section 3.7.1.1
\subsubsection[Duplication and Abbreviation]{\origsecnum{3.7.1.1}Duplication and Abbreviation}
\label{sec:chapters-ch03:3.7.1.1}

With the exception of minor variations in layout and epigraphy, inscriptions I,
2, 5, 8, 9, and 10 on the front of the shells and 1 through 8 on the back were ``carbon
copies,'' reproduced without abbreviation on all five plastrons.This conclusion depends, in several cases, on supplying graphs that would, it is assumed, have appeared on missing shell fragments. Due to the roughness of the inner surface of the shell, it is not always easy to see the graphs in the rubbings of the plastron backs; in these cases I rely on the k'ao-shih.83 The other inscriptions

% source: scan p. 99, section 3.7.3
\subsection[Context Method]{\origsecnum{3.7.3}Context Method}
\label{sec:chapters-ch03:3.7.3}

There are three crack notations on the front of the plastron pieces now remaining
to us. The first is a hsiao-chi, ``slightly auspicious,'' carved in the lower angle of the
front-inscription crack formed in hollow H6 of Ping-pien 14/15; I have associated this (in
table 11) with charge 8: ``Praying to lead away this sick tooth, the ting-sacrifice will be
favorable.'' The second crack notation is a shang-chi, ``highly auspicious,'' near the
top left corner of the right hypoplastron of Ping-pien 14; this notation lies in the lower

% End of extracted sections